% ============================================================
% 02-bezier-sympy-code-explained.tex
% Compile with:  pdflatex -shell-escape 02-bezier-sympy-code-explained.tex
% Run twice for table of contents.
% ============================================================

\documentclass[12pt, a4paper]{article}

% ---- Encoding & fonts ----------------------------------------
\usepackage[T1]{fontenc}
\usepackage[utf8]{inputenc}
\usepackage{lmodern}

% ---- Page geometry -------------------------------------------
\usepackage[
  top=2.5cm, bottom=2.5cm,
  left=2.8cm, right=2.8cm
]{geometry}

% ---- Colours -------------------------------------------------
\usepackage[dvipsnames]{xcolor}
\definecolor{codebg}{RGB}{248,248,248}
\definecolor{rulecolor}{RGB}{180,180,180}
\definecolor{examplebg}{RGB}{240,248,255}
\definecolor{examplerule}{RGB}{100,149,237}
\definecolor{notebg}{RGB}{255,250,230}
\definecolor{noterule}{RGB}{210,180,50}
\definecolor{diffbg}{RGB}{240,255,240}
\definecolor{diffrule}{RGB}{80,160,80}

% ---- Minted --------------------------------------------------
\usepackage{minted}
\setminted[julia]{
  bgcolor    = codebg,
  frame      = leftline,
  framerule  = 1.5pt,
  rulecolor  = rulecolor,
  linenos    = false,
  fontsize   = \small,
  breaklines = true,
  tabsize    = 4,
  style      = friendly,
}
\setminted[text]{
  bgcolor    = codebg,
  fontsize   = \small,
  breaklines = true,
}

% ---- Math ----------------------------------------------------
\usepackage{amsmath}
\usepackage{amssymb}

% ---- TikZ ----------------------------------------------------
\usepackage{tikz}
\usetikzlibrary{shapes.geometric, arrows.meta, positioning, fit, backgrounds}

% ---- Boxes ---------------------------------------------------
\usepackage{mdframed}

\newmdenv[
  backgroundcolor  = examplebg,
  linecolor        = examplerule, linewidth = 1.5pt,
  leftline=true, rightline=false, topline=false, bottomline=false,
  innerleftmargin=10pt, innerrightmargin=8pt,
  innertopmargin=6pt, innerbottommargin=6pt,
  skipabove=6pt, skipbelow=4pt,
]{examplebox}

\newmdenv[
  backgroundcolor  = notebg,
  linecolor        = noterule, linewidth = 1.5pt,
  leftline=true, rightline=false, topline=false, bottomline=false,
  innerleftmargin=10pt, innerrightmargin=8pt,
  innertopmargin=6pt, innerbottommargin=6pt,
  skipabove=6pt, skipbelow=4pt,
]{notebox}

\newmdenv[
  backgroundcolor  = diffbg,
  linecolor        = diffrule, linewidth = 1.5pt,
  leftline=true, rightline=false, topline=false, bottomline=false,
  innerleftmargin=10pt, innerrightmargin=8pt,
  innertopmargin=6pt, innerbottommargin=6pt,
  skipabove=6pt, skipbelow=4pt,
]{diffbox}

% ---- Hyperlinks ----------------------------------------------
\usepackage[colorlinks=true, linkcolor=black, urlcolor=MidnightBlue]{hyperref}

% ---- Misc ----------------------------------------------------
\usepackage{booktabs}
\usepackage{needspace}
\usepackage{parskip}
\usepackage{titlesec}
\usepackage{fancyhdr}
\usepackage{datetime2}

% ---- Section formatting --------------------------------------
\titleformat{\section}
  {\large\bfseries\color{MidnightBlue}}
  {\thesection.}{0.6em}{}[\vspace{-4pt}\rule{\linewidth}{0.4pt}]
\titleformat{\subsection}
  {\normalsize\bfseries\color{BrickRed}}
  {\thesubsection.}{0.5em}{}
\titleformat{\subsubsection}
  {\normalsize\bfseries}
  {\thesubsubsection.}{0.5em}{}

% ---- Header / footer -----------------------------------------
\pagestyle{fancy}
\fancyhf{}
\rhead{\textcolor{gray}{\small 02-bezier-sympy-code-explained}}
\lhead{\textcolor{gray}{\small B\'ezier Symbolic --- Code Explanation}}
\cfoot{\thepage}
\renewcommand{\headrulewidth}{0.3pt}

% ---- Helpers -------------------------------------------------
\newcommand{\cd}[1]{\texttt{#1}}
\newcommand{\NEW}{\textcolor{diffrule}{\textbf{[New]}}\;}

% ---- Title ---------------------------------------------------
\title{%
  \vspace{1cm}
  {\LARGE\bfseries B\'ezier Curve --- Algebraic Expansion (Symbolic)}\\[0.4em]
  {\large Code Explanation for Julia Beginners}\\[0.4em]
  {\normalsize\texttt{02-bezier-sympy.jl}}
}
\author{E.R.\ Nurwijayadi}
\date{\DTMtoday}


% =============================================================
\begin{document}
% =============================================================

\maketitle
\thispagestyle{empty}

\begin{center}
\textit{Directed and curated by E.R.\ Nurwijayadi.
Code and explanations AI-generated.}
\end{center}

\begin{center}
\begin{minipage}{0.85\textwidth}
\small
This document explains the Julia code in \cd{02-bezier-sympy.jl}
for readers new to Julia.
The script computes the same B\'ezier curve as \cd{01-bezier-algebra.jl}
but replaces the hand-written \cd{Poly} type with
\textbf{Symbolics.jl} --- Julia's computer algebra system.
This document is fully self-contained:
all concepts are explained from scratch even where they overlap
with the first script.
\end{minipage}
\end{center}

\vspace{1em}\hrule\vspace{1.5em}
\tableofcontents
\newpage


% =============================================================
\needspace{8\baselineskip}
\section{Type Map: How the Pieces Fit Together}
% =============================================================

\textit{Script~02 has a simpler structure than Script~01.
The hand-rolled \cd{Poly} type and all its operator overrides are gone.
In their place stands a single call to \cd{expand()} from Symbolics.jl,
which does the same work in one line.
The diagram below shows what remains.}

\bigskip

\tikzset{
  structbox/.style={
    rectangle, rounded corners=4pt,
    draw=MidnightBlue, line width=1.5pt,
    fill=examplebg,
    text width=3.2cm, align=center,
    font=\small\bfseries, inner sep=8pt
  },
  funcbox/.style={
    rectangle, rounded corners=2pt,
    draw=gray!60, line width=1pt,
    fill=white,
    text width=3.4cm, align=center,
    font=\small, inner sep=7pt
  },
  casbox/.style={
    rectangle, rounded corners=4pt,
    draw=diffrule, line width=1.5pt,
    fill=diffbg,
    text width=3.2cm, align=center,
    font=\small\bfseries, inner sep=8pt
  },
  entrybox/.style={
    rectangle, rounded corners=4pt,
    draw=BrickRed, line width=1.5pt,
    fill=notebg,
    text width=2.4cm, align=center,
    font=\small\bfseries, inner sep=8pt
  },
  myarrow/.style={-Stealth, thick, draw=gray!70},
  bluearrow/.style={-Stealth, thick, draw=MidnightBlue},
  redarrow/.style={-Stealth, thick, draw=BrickRed!80},
  arrowlabel/.style={font=\footnotesize\itshape, fill=white, inner sep=2pt}
}

\begin{center}
\begin{tikzpicture}[node distance=1.4cm and 2.6cm]

  % -- Symbolics.jl (external CAS, top left) -----------------
  \node[casbox] (symjl) {%
    \texttt{Symbolics.jl}\\[3pt]
    \footnotesize\normalfont
    \textit{external CAS library}\\
    \texttt{@variables u}\\
    \texttt{expand()}\\
    \texttt{substitute()}
  };

  % -- BezierExpansion struct (top right) --------------------
  \node[structbox, right=3.6cm of symjl] (bezier) {%
    \texttt{BezierExpansion}\\[3pt]
    \footnotesize\normalfont
    \texttt{basis::Vector}\\
    \texttt{exprs::Dict\{String,Any\}}\\
    \texttt{points, axes, n, degree, dim}
  };

  % -- BezierExpansion constructor + evaluate (below bezier) -
  \node[funcbox, below=1.6cm of bezier] (bezierctr) {%
    \texttt{BezierExpansion()}\\
    \texttt{evaluate()}\\[2pt]
    \footnotesize\normalfont
    \textit{constructor + evaluator}
  };

  % -- to_exact_rational (below symjl) -----------------------
  \node[funcbox, below=1.6cm of symjl] (ter) {%
    \texttt{to\_exact\_rational()}\\[2pt]
    \footnotesize\normalfont
    \textit{Symbolics result}\\
    \textit{to \texttt{Rational\{Int\}}}
  };

  % -- Formatting helpers (below ter) ------------------------
  \node[funcbox, below=1.2cm of ter] (fmt) {%
    \textbf{Formatting helpers}\\[2pt]
    \footnotesize\normalfont
    \texttt{fmt\_rational}\\
    \texttt{fmt\_sym}\\
    \texttt{fmt\_bernstein\_compact}
  };

  % -- Print functions (below bezierctr) ---------------------
  \node[funcbox, below=1.2cm of bezierctr] (print) {%
    \textbf{Print functions}\\[2pt]
    \footnotesize\normalfont
    \texttt{print\_step1} \ldots \texttt{print\_step5}\\
    \texttt{print\_all}
  };

  % -- main (far right) --------------------------------------
  \node[entrybox, right=2.0cm of bezier] (main) {%
    \texttt{main()}\\[2pt]
    \footnotesize\normalfont\normalcolor
    \textit{entry point}
  };

  % -- Arrows ------------------------------------------------
  \draw[bluearrow] (symjl)     -- node[arrowlabel, above]{provides} (bezier);
  \draw[bluearrow] (bezierctr) -- node[arrowlabel, right]{builds}   (bezier);
  \draw[myarrow]   (bezierctr) -- node[arrowlabel, left]{calls}     (symjl);
  \draw[myarrow]   (ter)       -- node[arrowlabel, left]{reads}     (symjl);
  \draw[myarrow]   (fmt)       -- node[arrowlabel, right]{uses}     (ter);
  \draw[myarrow]   (print)     -- node[arrowlabel, right]{reads}    (bezierctr);
  \draw[myarrow]   (print)     -- node[arrowlabel, above]{calls}    (fmt);
  \draw[redarrow]  (main)      -- node[arrowlabel, above]{calls}    (bezier);
  \draw[redarrow]  (main)      |- node[arrowlabel, right, pos=0.8]{calls} (print);

\end{tikzpicture}
\end{center}

\bigskip

\begin{notebox}
\textbf{Key difference from Script~01.}
Script~01 had a \cd{Poly} struct with overridden operators
(\cd{*}, \cd{+}, \cd{(p)(u)}) sitting between the constructor
and the results.
Here, \textbf{Symbolics.jl} (green box) replaces all of that.
\cd{expand()} does what \cd{expand\_bernstein} did.
\cd{substitute()} does what \cd{(p)(u)} did.
The arrow from \cd{BezierExpansion()} directly to Symbolics.jl
shows that the CAS is the new computation engine.
\end{notebox}


% =============================================================
\needspace{8\baselineskip}
\section{Imports and the Symbolic Variable}
% =============================================================

\begin{minted}{julia}
using Symbolics
using Printf

@variables u
\end{minted}

\needspace{5\baselineskip}
\subsection{\texttt{using Symbolics}}

\cd{using Symbolics} loads the Symbolics.jl package ---
Julia's built-in computer algebra system.
It provides the ability to work with symbolic mathematical
expressions: polynomials, equations, and expressions
containing variables like \cd{u} that have no fixed numeric value.

\needspace{5\baselineskip}
\subsection{\texttt{@variables u} --- creating a symbolic variable}

\cd{@variables u} is a \textbf{macro} call.
In Julia, macros begin with \cd{@} and transform code
before it runs.
This particular macro, provided by Symbolics.jl,
declares \cd{u} as a \textbf{symbolic variable}.

After this line, \cd{u} is not a number.
It is a placeholder --- a name that represents an unknown value.
You can build expressions with it:

\begin{examplebox}
\begin{minted}{julia}
@variables u
expr = 3*u^2 - 6*u + 1
\end{minted}

\cd{expr} is now the symbolic expression $3u^2 - 6u + 1$.
Julia does not compute a number --- it holds the expression
as a symbolic object.

You can then substitute a value later:
\begin{minted}{julia}
substitute(expr, Dict(u => 1//2))
\end{minted}
This gives the exact rational result $3 \cdot \tfrac{1}{4} - 6 \cdot \tfrac{1}{2} + 1 = \tfrac{3}{4} - 3 + 1 = -\tfrac{5}{4}$.
\end{examplebox}

\begin{notebox}
\textbf{Why a macro and not a function?}

A regular function call like \cd{declare\_variable("u")} would
create a variable with that string name, but you could not then
write \cd{u} in your code and have Julia know you mean that variable.
The macro \cd{@variables u} rewrites the source code so that
the Julia identifier \cd{u} is bound to a symbolic object.
That is something only a macro can do --- it operates on the
code itself before the program runs.
\end{notebox}


% =============================================================
\needspace{8\baselineskip}
\section{Exact Arithmetic: \texttt{Rational\{Int\}}}
% =============================================================

\begin{minted}{julia}
const CASE_DATA = Dict(
    2 => (
        points   = Rational{Int}[0 0 1;
                                  0 4 1;
                                  4 0 1;
                                  4 4 1;
                                  5 4 1],
        u_values = Rational{Int}[0, 1//4, 1//2, 3//4, 1],
        axes     = ["x", "y", "z"],
    ),
)
\end{minted}

\needspace{5\baselineskip}
\subsection{What is \texttt{Rational\{Int\}}?}

\cd{Rational\{Int\}} is Julia's built-in type for exact fractions.
It stores a numerator and denominator as plain integers,
with no floating-point representation involved.

The \cd{//} operator creates a \cd{Rational\{Int\}} value:

\begin{examplebox}
\begin{tabular}{lll}
\toprule
Expression & Value & Type \\
\midrule
\cd{1//4}  & $\frac{1}{4}$   & \cd{Rational\{Int\}} \\
\cd{3//4}  & $\frac{3}{4}$   & \cd{Rational\{Int\}} \\
\cd{4//1}  & $4$             & \cd{Rational\{Int\}} (same as integer 4, but exact) \\
\cd{1//3}  & $\frac{1}{3}$   & \cd{Rational\{Int\}} --- no rounding, ever \\
\bottomrule
\end{tabular}

\medskip
Compare: \cd{Float64(1/3) = 0.3333333333333333} (rounded).
\cd{1//3} is exact.
\end{examplebox}

\needspace{5\baselineskip}
\subsection{The typed array literal \texttt{Rational\{Int\}[...]}}

Writing \cd{Rational\{Int\}[0 0 1; 0 4 1; ...]} tells Julia
to create a matrix where every element is stored as
\cd{Rational\{Int\}}, even though the values look like plain integers.
So \cd{4} is stored as \cd{4//1}, not as \cd{Int64(4)} or
\cd{Float64(4.0)}.

This matters because all arithmetic that follows --- multiplying
coordinates by basis polynomials, summing terms --- will stay
in exact rational form throughout.
No floating-point error can creep in.

\needspace{5\baselineskip}
\subsection{Why exact arithmetic matters here}

The Bernstein basis coefficients are binomial coefficients
(always integers).
The control point coordinates are whole numbers stored as
\cd{Rational\{Int\}}.
The $u$ values are simple fractions: $0, \tfrac{1}{4}, \tfrac{1}{2},
\tfrac{3}{4}, 1$.

All of these are rational numbers.
Their polynomial combinations are therefore also rational numbers.
Script~02 \emph{computes} and \emph{displays} those exact fractions
--- for example $\frac{269}{256}$ --- rather than approximating
them as \cd{1.050781}.

\begin{diffbox}
\textbf{Compared to Script~01:}
Script~01 stored coordinates as \cd{Float64} and all arithmetic
was floating-point.
Script~02 stores coordinates as \cd{Rational\{Int\}} and all
arithmetic is exact integer fraction arithmetic.
The decimal output at the end is the same, but the
path to get there is fundamentally different.
\end{diffbox}


% =============================================================
\needspace{8\baselineskip}
\section{The \texttt{BezierExpansion} Struct}
% =============================================================

\begin{minted}{julia}
struct BezierExpansion
    points  ::Matrix{Rational{Int}}
    axes    ::Vector{String}
    n       ::Int
    degree  ::Int
    dim     ::Int
    basis   ::Vector        # symbolic expressions per B(n,i,u)
    exprs   ::Dict{String,Any}  # collected symbolic expr per axis
end
\end{minted}

Compare this to Script~01's struct:

\begin{diffbox}
\textbf{What changed in the struct:}

\medskip
\begin{tabular}{lll}
\toprule
Field & Script~01 & Script~02 \\
\midrule
\cd{points} & \cd{Matrix\{Float64\}} & \cd{Matrix\{Rational\{Int\}\}} \\
\cd{basis}  & \cd{Vector\{Poly\}}    & \cd{Vector} (untyped --- holds Symbolics expressions) \\
\cd{polys}  & \cd{Dict\{String,Poly\}} & renamed to \cd{exprs::Dict\{String,Any\}} \\
\bottomrule
\end{tabular}

\medskip
The fields \cd{axes}, \cd{n}, \cd{degree}, \cd{dim} are identical.
\end{diffbox}

\needspace{5\baselineskip}
\subsection{Why \texttt{Vector} and \texttt{Dict\{String,Any\}}?}

In Script~01, \cd{basis} was \cd{Vector\{Poly\}} --- a vector
where every element is a \cd{Poly} object, a type we defined ourselves.

In Script~02, \cd{basis} is just \cd{Vector} with no type parameter,
and \cd{exprs} values are \cd{Any}.
This is because Symbolics.jl expressions have complex internal types
that are awkward to name explicitly.
Saying \cd{Any} tells Julia: ``I don't want to restrict the type here,
accept whatever Symbolics.jl returns.''

This is less strict than Script~01, but it is the practical
approach when working with a CAS whose expression types
are not simple to predict.


% =============================================================
\needspace{8\baselineskip}
\section{The \texttt{BezierExpansion} Constructor}
% =============================================================

\begin{minted}{julia}
function BezierExpansion(points::Matrix{Rational{Int}},
                         axes::Vector{String})
    n   = size(points, 1)
    deg = n - 1
    dim = length(axes)

    basis = [
        expand(binomial(deg, i) * u^i * (1 - u)^(deg - i))
        for i in 0:n-1
    ]

    exprs = Dict{String,Any}()
    for (dim_idx, axis) in enumerate(axes)
        expr = sum(basis[i+1] * points[i+1, dim_idx] for i in 0:n-1)
        exprs[axis] = expand(expr)
    end

    BezierExpansion(points, axes, n, deg, dim, basis, exprs)
end
\end{minted}

This constructor does in about ten lines what Script~01 did
with \cd{expand\_bernstein}, \cd{Poly +}, \cd{Poly *}, and the
\cd{BezierExpansion} constructor combined.
The CAS absorbs all of that complexity.

\needspace{5\baselineskip}
\subsection{Building the basis: \texttt{expand()}}

\begin{minted}{julia}
basis = [
    expand(binomial(deg, i) * u^i * (1 - u)^(deg - i))
    for i in 0:n-1
]
\end{minted}

This is a \textbf{list comprehension} --- it evaluates the
expression in brackets for each value of \cd{i} from \cd{0}
to \cd{n-1} and collects the results into a vector.

For each \cd{i}, it builds the Bernstein polynomial
$B(n, i, u) = \binom{n}{i}\,u^i\,(1-u)^{n-i}$ symbolically,
then calls \cd{expand()} to distribute and collect all terms.

\begin{examplebox}
\textbf{For $n=4$ (quartic), $i=2$:}

\begin{minted}{julia}
expand(binomial(4, 2) * u^2 * (1 - u)^2)
\end{minted}

\cd{binomial(4, 2) = 6} (Julia built-in, exact integer)

The expression before expanding: $6 \cdot u^2 \cdot (1-u)^2$

After \cd{expand()}: $6u^2 - 12u^3 + 6u^4$

Symbolics.jl returns this as a symbolic expression object.
In Script~01, the same result came from \cd{expand\_bernstein(4, 2)},
which did the binomial expansion manually step by step.
Here, one call to \cd{expand()} handles everything.
\end{examplebox}

\needspace{5\baselineskip}
\subsection{Building per-axis polynomials}

\begin{minted}{julia}
exprs = Dict{String,Any}()
for (dim_idx, axis) in enumerate(axes)
    expr = sum(basis[i+1] * points[i+1, dim_idx] for i in 0:n-1)
    exprs[axis] = expand(expr)
end
\end{minted}

For each axis, it computes the weighted sum of all basis polynomials,
each multiplied by the corresponding coordinate.

\cd{sum(...for i in 0:n-1)} is Julia's generator sum syntax:
it accumulates $\sum_{i=0}^{n-1} B_i \cdot \text{coord}_i$
without building an intermediate array.

The final \cd{expand(expr)} collects all like powers of \cd{u}
into a clean monomial polynomial.

\begin{examplebox}
\textbf{For the $x$ axis, Case~2:}

$x$-coordinates are $0, 0, 4, 4, 5$.

\[
  \text{expr} = B_0 \cdot 0 + B_1 \cdot 0 + B_2 \cdot 4 + B_3 \cdot 4 + B_4 \cdot 5
\]

After \cd{expand()}:
\[
  x(u) = 24u^2 - 32u^3 + 13u^4
\]

Stored as \cd{exprs["x"]} --- a Symbolics expression,
not a \cd{Poly} object.
\end{examplebox}


% =============================================================
\needspace{8\baselineskip}
\section{The \texttt{evaluate} Function}
% =============================================================

\begin{minted}{julia}
function evaluate(exp::BezierExpansion, u_val::Rational{Int})
    [substitute(exp.exprs[ax], Dict(u => u_val)) for ax in exp.axes]
end
\end{minted}

\cd{substitute(expr, Dict(u => u\_val))} is the Symbolics.jl
way of evaluating a symbolic expression at a specific value.
It replaces every occurrence of the symbolic variable \cd{u}
with the given rational value \cd{u\_val}.

The \cd{Dict(u => u\_val)} syntax creates a substitution map:
``replace \cd{u} with \cd{u\_val}.''

\begin{examplebox}
\textbf{Evaluating $x(u)$ at $u = \tfrac{1}{2}$:}

\begin{minted}{julia}
substitute(exprs["x"], Dict(u => 1//2))
\end{minted}

This replaces $u$ with $\tfrac{1}{2}$ in $24u^2 - 32u^3 + 13u^4$:
\[
  24 \cdot \tfrac{1}{4} - 32 \cdot \tfrac{1}{8} + 13 \cdot \tfrac{1}{16}
  = 6 - 4 + \tfrac{13}{16} = \tfrac{45}{16}
\]

The result is a Symbolics expression representing $\tfrac{45}{16}$.
It is not yet a plain Julia \cd{Rational\{Int\}} ---
that conversion happens in \cd{to\_exact\_rational()}.
\end{examplebox}

\begin{diffbox}
\textbf{Compared to Script~01:}
Script~01 evaluated a polynomial by calling \cd{(p)(u\_val)},
which used Horner's method on the \cd{Float64} coefficient vector.
Script~02 calls \cd{substitute()}, which is symbolic substitution
--- the CAS replaces the variable and simplifies algebraically.
The result is an exact rational, not a float.
\end{diffbox}


% =============================================================
\needspace{8\baselineskip}
\section{Formatting Helpers}
% =============================================================

\needspace{5\baselineskip}
\subsection{\texttt{fmt\_rational} --- multiple dispatch on two types}

\begin{minted}{julia}
fmt_rational(r::Rational{Int}) =
    isone(denominator(r)) ? string(numerator(r)) : string(r)
fmt_rational(n::Integer) = string(n)
\end{minted}

This is a compact example of \textbf{multiple dispatch}:
two methods with the same name but different argument types.
Julia picks the right one based on what you pass in.

\textbf{First method} handles a \cd{Rational\{Int\}}:
\begin{itemize}
  \item \cd{denominator(r)} returns the bottom of the fraction.
  \item \cd{isone(denominator(r))} checks if it equals 1.
  \item If the denominator is 1, the fraction is a whole number ---
        print just the numerator (e.g.\ \cd{"4"} not \cd{"4//1"}).
  \item Otherwise print the full fraction (e.g.\ \cd{"269//256"}).
\end{itemize}

\textbf{Second method} handles any plain \cd{Integer}
(like \cd{Int64}): just convert to string directly.

\begin{examplebox}
\cd{fmt\_rational(4//1)}    $\to$ \cd{"4"} \quad (denominator is 1, so whole number) \\
\cd{fmt\_rational(269//256)} $\to$ \cd{"269//256"} \quad (denominator is 256) \\
\cd{fmt\_rational(5)}       $\to$ \cd{"5"} \quad (plain integer, second method)
\end{examplebox}

\needspace{5\baselineskip}
\subsection{\texttt{fmt\_sym} --- formatting a Symbolics expression}

\begin{minted}{julia}
fmt_sym(expr) = string(expr)
\end{minted}

This one-liner converts any Symbolics.jl expression to a
human-readable string.
It has no type annotation on \cd{expr} ---
it accepts anything and calls Julia's built-in \cd{string()}
on it, which Symbolics.jl has extended to produce readable output.

This is how \cd{(24//1)*(u\^{}2) - (32//1)*(u\^{}3) + (13//1)*(u\^{}4)}
ends up in the printed output.

\needspace{5\baselineskip}
\subsection{\texttt{to\_exact\_rational} --- converting a CAS result to a fraction}

\begin{minted}{julia}
function to_exact_rational(val)
    v = Symbolics.value(val)
    try
        r = rationalize(Int, Float64(v); tol=1e-12)
        return r
    catch
        return v
    end
end
\end{minted}

After \cd{substitute()} returns a Symbolics expression like
\cd{45//16}, this function extracts the plain Julia
\cd{Rational\{Int\}} value from it.
This is needed because Symbolics.jl wraps even simple
numeric results in its own expression type, and we want
a plain fraction we can format and convert to \cd{Float64}.

The steps are:
\begin{enumerate}
  \item \cd{Symbolics.value(val)} --- strips the Symbolics wrapper,
        returning a lower-level numeric type.
  \item \cd{Float64(v)} --- converts to a float temporarily.
  \item \cd{rationalize(Int, ...; tol=1e-12)} --- converts the
        float back to the nearest \cd{Rational\{Int\}} within
        a tolerance of $10^{-12}$. This recovers the exact
        fraction, since the float was exact to begin with
        (it came from rational arithmetic).
  \item The \cd{try/catch} handles any edge cases where the
        conversion fails --- in that case the raw value \cd{v}
        is returned as-is.
\end{enumerate}

\begin{examplebox}
\textbf{Tracing \cd{to\_exact\_rational} for $x\!\left(\tfrac{1}{2}\right)$:}

\medskip
After \cd{substitute}, \cd{val} is a Symbolics expression
representing $\tfrac{45}{16}$.

\cd{Symbolics.value(val)} $\to$ a SymbolicUtils numeric for $2.8125$

\cd{Float64(v)} $\to$ \cd{2.8125}

\cd{rationalize(Int, 2.8125; tol=1e-12)} $\to$ \cd{45//16}

Result: \cd{45//16} as a plain Julia \cd{Rational\{Int\}}.
\end{examplebox}

\begin{notebox}
\textbf{Why go through Float64 and back?}

It seems roundabout to convert a rational to float and then
back to rational.
The reason is that \cd{Symbolics.value()} does not always
return a type with a direct \cd{Rational} conversion path ---
going through \cd{Float64} and \cd{rationalize} is the
reliable general approach.
For the values in this program (small integer fractions),
the round-trip is exact.
\end{notebox}

\needspace{5\baselineskip}
\subsection{\texttt{fmt\_bernstein\_compact}}

\begin{minted}{julia}
const SUPERSCRIPTS = Dict(0=>"", 1=>"", 2=>"^2", 3=>"^3",
                          4=>"^4", 5=>"^5")

function fmt_bernstein_compact(n::Int, i::Int)
    b      = binomial(n, i)
    m      = n - i
    u_part = i == 0 ? "" : (i == 1 ? "u" : "u$(get(SUPERSCRIPTS,i,"^$i"))")
    v_part = m == 0 ? "" : (m == 1 ? "(1-u)" :
                            "(1-u)$(get(SUPERSCRIPTS,m,"^$m"))")
    b_str  = b == 1 ? "" : string(b)
    result = "$(b_str)$(u_part)$(v_part)"
    isempty(result) ? "1" : result
end
\end{minted}

This produces compact human-readable symbolic strings like
\cd{"6u\^{}2(1-u)\^{}2"} for the print steps.
It is structurally identical to \cd{fmt\_bernstein\_symbolic}
in Script~01, with one extra guard at the end:
\cd{isempty(result) ? "1" : result}.
This handles the edge case $B(0,0,u) = 1$, which would
otherwise produce an empty string.

\begin{examplebox}
\begin{tabular}{lll}
\toprule
Call & Formula & Output \\
\midrule
\cd{fmt\_bernstein\_compact(4, 0)} & $(1-u)^4$        & \cd{"(1-u)\^{}4"} \\
\cd{fmt\_bernstein\_compact(4, 1)} & $4u(1-u)^3$      & \cd{"4u(1-u)\^{}3"} \\
\cd{fmt\_bernstein\_compact(4, 2)} & $6u^2(1-u)^2$    & \cd{"6u\^{}2(1-u)\^{}2"} \\
\cd{fmt\_bernstein\_compact(4, 4)} & $u^4$            & \cd{"u\^{}4"} \\
\bottomrule
\end{tabular}
\end{examplebox}


% =============================================================
\needspace{8\baselineskip}
\section{Print Step Functions}
% =============================================================

\textit{The five print step functions in Script~02 do exactly
the same job as in Script~01 --- they display the five stages
of the algebraic expansion.
The structure of each function is identical.
What changed is which fields they read from \cd{exp}
and which formatting helpers they call.}

\needspace{5\baselineskip}
\subsection{What is the same}

All five functions follow the same pattern:
call \cd{print\_step\_header}, loop over axes and/or control
points, call a formatting helper, call \cd{println}.
None of them return a value.

\begin{center}
\begin{tabular}{lp{5cm}p{5cm}}
\toprule
Step & Script~01 source & Script~02 source \\
\midrule
Step~1 & \cd{fmt\_bernstein\_symbolic} & \cd{fmt\_bernstein\_compact} \\
Step~2 & \cd{fmt\_bernstein\_symbolic} + \cd{fmt\_coeff} & \cd{fmt\_bernstein\_compact} + \cd{fmt\_rational} \\
Step~3 & \cd{fmt\_poly(exp.basis[i+1])} & \cd{fmt\_sym(exp.basis[i+1])} \\
Step~4 & \cd{expand(basis[i]*coord)} via \cd{Poly*} & \cd{expand(exp.basis[i]*coord)} via Symbolics \\
Step~5 & \cd{evaluate} returns \cd{Float64} & \cd{evaluate} returns Symbolics expr, then \cd{to\_exact\_rational} \\
\bottomrule
\end{tabular}
\end{center}

\needspace{5\baselineskip}
\subsection{Step 5: the new evaluation loop}

The evaluation loop in Step~5 is the most changed part:

\begin{minted}{julia}
for u_val in u_values
    results = evaluate(exp, u_val)
    row = "  " * lpad(fmt_rational(u_val), col_u)
    for val in results
        r     = to_exact_rational(val)
        exact = fmt_rational(r)
        dec   = @sprintf("%.6f", Float64(r))
        row  *= "  " * lpad(exact, col_exact) *
                "  " * lpad(dec, col_dec)
    end
    println(row)
end
\end{minted}

For each \cd{u\_val}, it calls \cd{evaluate} (which calls
\cd{substitute} internally), then for each result:

\begin{enumerate}
  \item \cd{to\_exact\_rational(val)} --- extract the exact fraction
  \item \cd{fmt\_rational(r)} --- format as \cd{"45//16"} or \cd{"5"}
  \item \cd{Float64(r)} --- convert to decimal for the second column
\end{enumerate}

This produces the two-column output (exact + decimal) that
distinguishes Script~02's evaluation table from Script~01's.

\begin{examplebox}
\textbf{One row traced: $u = \tfrac{1}{4}$}

\medskip
\cd{fmt\_rational(1//4)} $\to$ \cd{"1//4"} (printed in the $u$ column)

For the $x$ axis:
\begin{enumerate}
  \item \cd{evaluate} calls \cd{substitute(exprs["x"], Dict(u => 1//4))}
        $\to$ Symbolics result for $\tfrac{269}{256}$
  \item \cd{to\_exact\_rational} extracts \cd{269//256}
  \item \cd{fmt\_rational(269//256)} $\to$ \cd{"269//256"}
  \item \cd{Float64(269//256)} $\to$ \cd{1.05078125} $\to$ \cd{"1.050781"}
\end{enumerate}
Row printed: \cd{1//4 \quad 269//256 \quad 1.050781 \quad ...}
\end{examplebox}


% =============================================================
\needspace{8\baselineskip}
\section{Summary: What Symbolics.jl Replaced}
% =============================================================

\begin{center}
\begin{tabular}{p{5cm}p{5cm}}
\toprule
Script~01 (hand-written) & Script~02 (Symbolics.jl) \\
\midrule
\cd{struct Poly} with \cd{coeffs::Vector\{Float64\}} &
  Symbolics expression (no separate type) \\[4pt]
\cd{Poly * scalar} (operator override) &
  \cd{expr * coord} (Symbolics handles \cd{*}) \\[4pt]
\cd{Poly + Poly} (operator override, zero-padding) &
  \cd{expr + expr} (Symbolics handles \cd{+}) \\[4pt]
\cd{expand\_bernstein(n, i)} (manual binomial steps) &
  \cd{expand(binomial(n,i) * u\^{}i * (1-u)\^{}m)} \\[4pt]
\cd{(p)(u)} functor via Horner's method &
  \cd{substitute(expr, Dict(u => val))} \\[4pt]
\cd{Float64} results &
  \cd{Rational\{Int\}} exact fractions \\[4pt]
One column per axis in evaluation table &
  Two columns: exact fraction + decimal \\
\bottomrule
\end{tabular}
\end{center}

\begin{notebox}
\textbf{Which is better?}

Neither --- they make different trade-offs.

Script~01 is self-contained: it needs no external library,
every operation is explicit and transparent, and it is easy
to follow step by step.

Script~02 is more powerful: it produces exact results,
handles any rational input without rounding, and the
constructor is dramatically shorter because the CAS
absorbs all the polynomial algebra.

\textit{Script~01 teaches you how the engine works.
Script~02 lets you use the engine without building it.
Both are worth understanding.}
\end{notebox}

\vspace{2em}\hrule\vspace{0.5em}
{\small\textcolor{gray}{%
Code explanation for \texttt{02-bezier-sympy.jl}
\textbar{} Directed and curated by E.R.\ Nurwijayadi
\textbar{} Code and explanations AI-generated}}

\end{document}
